% Reviewed translation: PDF pages 88--92 / printed pages 74--78.
% Stops on PDF page 92 immediately before Section 5 begins.

\MathBookSourcePage{88}
所以收敛关系 
\eqref{eq:abstract-dual-gradient-convergence} 蕴含
\[
\lim_{m\to\infty}\lambda^m=\lambda
\quad\text{in }M,
\]
再在 
\eqref{eq:abstract-um-from-lambda} 中取极限, 即得到 $u^m$ 的收敛性.
\end{Proof}

\setcounter{statement}{5}
\begin{Remark}\label{rem:abstract-descent-angle-meaning}
收敛判据 
\eqref{eq:abstract-descent-angle-condition} 蕴含下降方向 $\omega^m$
永远不会渐近地正交于梯度 $g^m$.
\end{Remark}

作为第一个应用, 沿梯度下降便得到带最优参数的简单梯度算法:
\begin{equation}\label{eq:abstract-simple-gradient-direction}
\omega^m=g^m.
\end{equation}
所以这个算法的一般步骤是: 对 $m\ge0$, 已知
$(u^m,\lambda^m)\in X\times M$, 由下式确定
$(z^m,g^m)\in X\times M$, $\rho^m\in\mathbb R$ 以及
$(u^{m+1},\lambda^{m+1})\in X\times M$:
\begin{equation}\label{eq:abstract-simple-gradient-system}
\left\{
\begin{aligned}
Cg^m&=\chi-Bu^m,\\
A_rz^m&=B'g^m,\\
\rho^m&=c(g^m,g^m)/b(z^m,g^m),\\
\lambda^{m+1}&=\lambda^m-\rho^mg^m,\\
u^{m+1}&=u^m+\rho^mz^m.
\end{aligned}
\right.
\end{equation}
这里判据 
\eqref{eq:abstract-descent-angle-condition} 显然以 $\alpha=1$ 成立.
于是有:

\setcounter{statement}{3}
\begin{Corollary}\label{cor:abstract-optimal-simple-gradient-convergence}
设 $b(\cdot,\cdot)$, $a_r(\cdot,\cdot)$ 和 $c(\cdot,\cdot)$
满足 Theorem~\ref{thm:abstract-general-descent-convergence} 的假设.
则带最优参数的简单梯度算法
\eqref{eq:abstract-simple-gradient-system} 收敛.
\end{Corollary}

这里值得指出, 即使参数 $\rho^m$ 不是最优的, 只要它们位于适当区间内,
简单梯度算法仍然收敛. 特别强调, 所得格式的收敛性不要求
$a(\cdot,\cdot)$ 对称. 其证明方法与
Theorem~\ref{thm:abstract-general-descent-convergence} 完全不同.

\setcounter{statement}{6}
\begin{Theorem}\label{thm:abstract-nonoptimal-gradient-convergence}
保留 Corollary~\ref{cor:abstract-optimal-simple-gradient-convergence} 的假设,
但不一定假设形式 $a(\cdot,\cdot)$ 对称. 则算法
\begin{equation}\tag{\ref*{eq:abstract-general-descent-lambda}$'$}
\label{eq:abstract-nonoptimal-lambda-update}
\lambda^{m+1}=\lambda^m-\rho^mg^m,
\end{equation}
以及
\begin{equation*}\tag{\ref*{eq:abstract-um-from-lambda}}
a_r(u^m,v)=\langle l,v\rangle
+r b(v,C^{-1}\chi)-b(v,\lambda^m)
\qquad\forall v\in X
\end{equation*}
对每个初值 $\lambda^0\in M$ 以及每个满足
\begin{equation}\label{eq:abstract-nonoptimal-step-range}
0<\inf_m\rho^m\le\sup_m\rho^m<2\gamma\widetilde\alpha_r
\end{equation}
的 $\rho^m$ 选择都收敛, 其中
\begin{equation}\tag{\ref*{eq:abstract-ar-ellipticity}$'$}
\label{eq:abstract-alpha-r-tilde}
\widetilde\alpha_r
=\inf_{v\in X}\left[a_r(v,v)/\|Bv\|_{M'}^2\right].
\end{equation}
\end{Theorem}

\MathBookSourcePage{89}
\begin{Proof}
置
\[
e^m=u^m-u,
\qquad
v^m=\lambda^m-\lambda.
\]
由 
\eqref{eq:abstract-q-first},
\eqref{eq:abstract-q-second} 和
\eqref{eq:abstract-um-from-lambda}, 得
\[
a_r(e^m,v)+b(v,v^m)=0
\qquad\forall v\in X,
\]
\[
-c(v^{m+1}-v^m,\mu)+\rho^mb(e^m,\mu)=0
\qquad\forall\mu\in M.
\]
在上述方程中取 $v=e^m$ 和 $\mu=v^{m+1}$, 得
\[
\begin{aligned}
c(v^{m+1}-v^m,v^{m+1})
&=\rho^m b(e^m,v^{m+1})\\
&=\rho^m\{b(e^m,v^{m+1}-v^m)-a_r(e^m,e^m)\}.
\end{aligned}
\]
因为双线性形式 $c(\cdot,\cdot)$ 对称, 可写成
\[
2c(v^{m+1}-v^m,v^{m+1})
=c(v^{m+1})-c(v^m)+c(v^{m+1}-v^m),
\]
其中 $c(\mu)$ 表示 $c(\mu,\mu)$. 因而利用
\eqref{eq:abstract-c-ellipticity}, 得
\[
\begin{aligned}
c(v^{m+1})-c(v^m)
+\gamma\|v^{m+1}-v^m\|_M^2
+2\rho^m a_r(e^m,e^m)\\
\le2\rho^m\|Be^m\|_{M'}
\|v^{m+1}-v^m\|_M.
\end{aligned}
\]
利用不等式 $2ab\le\gamma a^2+(1/\gamma)b^2$ 和
\eqref{eq:abstract-alpha-r-tilde}, 得
\[
2\rho^m\|Be^m\|_{M'}\|v^{m+1}-v^m\|_M
\le\gamma\|v^{m+1}-v^m\|_M^2
+(\rho^m)^2a_r(e^m,e^m)/(\widetilde\alpha_r\gamma).
\]
于是应用 
\eqref{eq:abstract-ar-ellipticity}, 得
\[
c(v^{m+1})-c(v^m)
+\rho^m\{2-\rho^m/(\widetilde\alpha_r\gamma)\}
\alpha_r\|e^m\|_X^2\le0.
\]
因为 $\rho^m$ 满足 
\eqref{eq:abstract-nonoptimal-step-range},
存在与 $m$ 无关的常数 $\delta>0$, 使
\[
2\widetilde\alpha_r-\rho^m/\gamma\ge\delta.
\]
所以
\begin{equation}\label{eq:abstract-nonoptimal-error-decay-step}
c(v^{m+1})-c(v^m)
+\rho^m\delta(\alpha_r/\widetilde\alpha_r)
\|e^m\|_X^2\le0.
\end{equation}
另一方面, inf-sup 条件 
\eqref{eq:abstract-inf-sup} 给出
\[
\begin{aligned}
\|v^m\|_M
&\le(1/\beta)\sup_{v\in X}
\frac{b(v,v^m)}{\|v\|_X}\\
&=(1/\beta)\sup_{v\in X}
\frac{a_r(e^m,v)}{\|v\|_X}
\le(\|a_r\|/\beta)\|e^m\|_X,
\end{aligned}
\]
其中 $\|a_r\|$ 表示形式 $a_r(\cdot,\cdot)$ 的范数. 因此
\[
\|e^m\|_X^2
\ge\beta^2c(v^m)/(\|c\|\|a_r\|^2),
\]
而 
\eqref{eq:abstract-nonoptimal-error-decay-step} 蕴含
\[
c(v^{m+1})\le c(v^m)
\{1-(\alpha_r/\widetilde\alpha_r)
(\rho^m\delta\beta^2)/(\|c\|\|a_r\|^2)\}.
\]

\MathBookSourcePage{90}
置 $\rho=\inf_m\rho^m$. 由
\eqref{eq:abstract-nonoptimal-step-range}, $\rho>0$. 因而
\begin{equation}\label{eq:abstract-linear-error-decay}
c(v^m)\le\theta^m c(v^0),
\qquad
0\le\theta=1-(\alpha_r/\widetilde\alpha_r)
(\rho\delta\beta^2)/(\|c\|\|a_r\|^2)<1.
\end{equation}
所以序列 $c(v^m)$ 收敛到零, 且
\[
\lim_{m\to\infty}\|\lambda^m-\lambda\|_M=0.
\]
进而这蕴含 $(u^m)$ 收敛.
\end{Proof}

\setcounter{statement}{6}
\begin{Remark}\label{rem:abstract-linear-convergence-rate}
由 
\eqref{eq:abstract-linear-error-decay} 可知, 每次迭代都至少以
$\theta^{1/2}<1$ 的因子缩减误差范数 $(c(v^m))^{1/2}$.
因此迭代方法
\eqref{eq:abstract-nonoptimal-lambda-update} 与
\eqref{eq:abstract-um-from-lambda} 线性收敛.
\end{Remark}

\setcounter{statement}{7}
\begin{Remark}\label{rem:abstract-constant-gradient-step}
用常参数 $\rho$ 代替 $\rho^m$ 的简单梯度算法也被广泛使用.
当然, 一次固定 $\rho$ 会减少计算量, 但也会降低方法的收敛速度.
理论和实际论证都引导我们按如下准则选择 $r$ 和 $\rho$:

\textup{(i)} $r$ 尽可能大, 因为 $A_r$ 的条件数随 $r$ 增大;

\textup{(ii)} 取
\[
\rho=(r\gamma^2)/\|c\|^2.
\]
事实上,
\[
\langle C^{-1}\chi,\chi\rangle
\ge(\gamma/\|c\|^2)\|\chi\|_{M'}^2.
\]
当 $a(\cdot,\cdot)$ 在 $X$ 上半正定(见
\eqref{eq:abstract-semipositive})时, 这蕴含
\[
a_r(v,v)\ge r\langle C^{-1}Bv,Bv\rangle
\ge(\gamma r/\|c\|^2)\|Bv\|_{M'}^2,
\]
从而
\[
\widetilde\alpha_r\ge\gamma r/\|c\|^2.
\]
所以以上 $\rho$ 满足 
\eqref{eq:abstract-nonoptimal-step-range}.
更具体地, 当 $M=M'$ 且 $C=\operatorname{Id}$ 时, 可以取 $\rho=r$.
\end{Remark}

现在转向共轭梯度法. 不再像 
\eqref{eq:abstract-simple-gradient-direction} 那样沿梯度方向
(即最速下降方向)更新 $\lambda^m$, 而是对超曲面
$K_r=\text{constant}$, 取与前一次下降方向共轭的方向. 这等价于取
\begin{equation}\label{eq:abstract-conjugate-direction}
\omega^m=g^m+\sigma^m\omega^{m-1},
\end{equation}
其中实参数 $\sigma^m$ 的选择使
\[
D^2K_r\mathbin{:}(\omega^m,\omega^{m-1})=0,
\]

\MathBookSourcePage{91}
即
\begin{equation}\label{eq:abstract-conjugate-sigma-definition}
\sigma^m=-\frac{D^2K_r\mathbin{:}(g^m,\omega^{m-1})}
{D^2K_r\mathbin{:}(\omega^{m-1},\omega^{m-1})}.
\end{equation}
事实上, $\rho^m$ 和 $\sigma^m$ 有更便于使用的表达式.
回顾如下经典 Theorem(在每一本讨论共轭梯度法的著作中都能找到).

\setcounter{statement}{7}
\begin{Theorem}\label{thm:abstract-conjugate-orthogonality}
有如下正交关系:
\[
D^2K_r\mathbin{:}(\omega^i,\omega^j)=0
\quad\text{if }i\ne j,
\]
\[
c(g^i,g^j)=0
\quad\text{if }i\ne j,
\]
\[
c(g^i,\omega^j)=0
\quad\text{if }i>j.
\]
\end{Theorem}

因此可以简化 
\eqref{eq:abstract-optimal-step} 和
\eqref{eq:abstract-conjugate-sigma-definition}.

\setcounter{statement}{4}
\begin{Corollary}\label{cor:abstract-conjugate-parameters}
参数 $\rho^m$ 和 $\sigma^m$ 满足
\begin{equation}\label{eq:abstract-conjugate-parameter-formulas}
\left\{
\begin{aligned}
\sigma^m&=c(g^m,g^m)/c(g^{m-1},g^{m-1}),\\
\rho^m&=c(g^m,g^m)/
\{D^2K_r\mathbin{:}(\omega^m,g^m)\}.
\end{aligned}
\right.
\end{equation}
\end{Corollary}
\begin{Proof}
先回顾
\[
c(g^m,v)=DK_r(\lambda^m)\mathbin{:}v
\qquad\forall v\in M.
\]
其次, 仿照 Theorem~\ref{thm:abstract-general-descent-convergence}, 得
\[
DK_r(\lambda^{m+1})-DK_r(\lambda^m)
=D^2K_r(\lambda^{m+1}-\lambda^m)
=-\rho^mD^2K_r\omega^m,
\]
即
\begin{equation}\label{eq:abstract-gradient-difference}
c(g^{m+1}-g^m,v)
=-\rho^mD^2K_r\mathbin{:}(\omega^m,v)
\qquad\forall v\in M.
\end{equation}
另一方面, 因为 $D^2K_r$ 自伴(见
\eqref{eq:abstract-kr-second-derivative}), 有
\[
\begin{aligned}
\sigma^m
&=-\frac{\rho^{m-1}D^2K_r\mathbin{:}(\omega^{m-1},g^m)}
{\rho^{m-1}D^2K_r\mathbin{:}(\omega^{m-1},\omega^{m-1})}\\
&=-\frac{c(g^m-g^{m-1},g^m)}
{c(g^m-g^{m-1},\omega^{m-1})}.
\end{aligned}
\]
于是 Theorem~\ref{thm:abstract-conjugate-orthogonality} 和
\eqref{eq:abstract-conjugate-direction} 给出
\[
\begin{aligned}
\sigma^m
&=c(g^m,g^m)/c(g^{m-1},\omega^{m-1})\\
&=c(g^m,g^m)/c(g^{m-1},g^{m-1}).
\end{aligned}
\]
最后在 
\eqref{eq:abstract-gradient-difference} 中取 $v=g^m$,
立即得到 $\rho^m$ 所需的表达式.
\end{Proof}

利用这个 Corollary, 可以更高效地组织共轭梯度算法的一般步骤.
对 $m\ge0$, 已知 $(u^m,\lambda^m)\in X\times M$, 由下式计算
$g^m,\omega^m\in M$, $z^m\in X$, $\rho^m,\sigma^m\in\mathbb R$
以及 $(u^{m+1},\lambda^{m+1})\in X\times M$:

\MathBookSourcePage{92}
\begin{equation}\label{eq:abstract-conjugate-gradient-system}
\left\{
\begin{aligned}
Cg^m&=\chi-Bu^m,\\
\sigma^m&=c(g^m,g^m)/c(g^{m-1},g^{m-1})
&&\text{only if }m\ge1,\\
\omega^m&=g^m+\sigma^m\omega^{m-1}
&&\text{only if }m\ge1,\\
\omega^0&=g^0&&\text{otherwise},\\
A_rz^m&=B'\omega^m,\\
\rho^m&=c(g^m,g^m)/b(z^m,g^m),\\
\lambda^{m+1}&=\lambda^m-\rho^m\omega^m,\\
u^{m+1}&=u^m+\rho^mz^m.
\end{aligned}
\right.
\end{equation}
注意, 每次迭代的计算量略大于简单梯度算法,
但主要计算工作仍在 $A_r$ 的求逆上.
最后, 容易证明这个格式收敛.

\setcounter{statement}{5}
\begin{Corollary}\label{cor:abstract-conjugate-gradient-convergence}
在 Corollary~\ref{cor:abstract-optimal-simple-gradient-convergence} 的假设下,
共轭梯度算法 
\eqref{eq:abstract-conjugate-gradient-system} 收敛.
\end{Corollary}
\begin{Proof}
注意, 收敛判据 
\eqref{eq:abstract-descent-angle-condition} 等价于
\begin{equation}\label{eq:abstract-conjugate-step-criterion}
\rho^m\ge\alpha'\|g^m\|_M/\|\omega^m\|_M
\qquad\text{for some constant }\alpha'>0.
\end{equation}
事实上, 由 
\eqref{eq:abstract-general-descent-system} 和 $z^m$ 的表达式,
\[
\rho^m
=c(g^m,\omega^m)/b(z^m,\omega^m)
=c(g^m,\omega^m)/
\{D^2K_r\mathbin{:}(\omega^m,\omega^m)\}.
\]
因此 
\eqref{eq:abstract-descent-angle-condition} 和
\eqref{eq:abstract-kr-hessian-bounds} 蕴含
\eqref{eq:abstract-conjugate-step-criterion}, 其中 $\alpha'=\alpha/\tau_r$;
反之, 
\eqref{eq:abstract-conjugate-step-criterion} 和
\eqref{eq:abstract-kr-hessian-bounds} 蕴含
\eqref{eq:abstract-descent-angle-condition}, 其中 $\alpha=\alpha'\delta_r$.

其次注意
\[
D^2K_r\mathbin{:}(\omega^m,g^m)
=D^2K_r\mathbin{:}(\omega^m,\omega^m)>0.
\]
因而 
\eqref{eq:abstract-conjugate-step-criterion} 是
\eqref{eq:abstract-conjugate-parameter-formulas} 的直接推论.
所以该格式收敛.
\end{Proof}
% Section 5 begins later on the same PDF page 92 and is not included here.
